\section{The PUT Base Model}
\label{sec:put-base-model}

The base privacy-utility trade-off (PUT) architecture this paper
studies, following Mo-PAE~\cite{zhan2023mopae}, has three
components: a Data Reconstruction Unit (DRU), a Mobility Prediction
Unit (MPU), and a User Re-identification Unit (URU). Raw mobility data
$\mathcal{X}$ is passed through an encoder $\mathrm{Enc}_L$ to learn a
representation $\mathcal{F} = \mathrm{Enc}_L(\mathcal{X}; \theta_E)$.
A decoder $\mathrm{Dec}_L$ attempts to reconstruct $\mathcal{X}$ from
$\mathcal{F}$ (the DRU, trained as a built-in adversary minimizing
reconstruction loss $\mathcal{L}_R$); the URU
$P_D(\mathcal{F};\theta_P')$ outputs a probability distribution over
user identities and is trained to minimize re-identification loss
$\mathcal{L}_P$; the MPU $U_D(\mathcal{F};\theta_U')$ outputs a
probability distribution over the next location and is trained to
minimize prediction loss $\mathcal{L}_U$. The encoder is trained
adversarially against the URU and DRU while cooperating with the MPU:
it must retain information useful for next-location prediction
(minimizing $\mathcal{L}_U$) while destroying information useful for
reconstruction and re-identification (maximizing $\mathcal{L}_R$ and
$\mathcal{L}_P$), giving the min-max training objective
\begin{equation}
\min\mathcal{L}_{\mathrm{sum}} = \min_{\mathcal{L}_U}\ \max_{\mathcal{L}_R,\mathcal{L}_P}
\sum_{i} \big(\mathcal{L}_U(f_i),\ \mathcal{L}_P(f_i),\ \mathcal{L}_R(f_i)\big).
\label{eq:sumloss}
\end{equation}
A fairness term $\mathcal{L}_F$ (Section~\ref{sec:individual-fairness})
can be added to Equation~\ref{eq:sumloss} with its own weight
$\lambda_4$, alongside the weights $\lambda_1,\lambda_2,\lambda_3$
already governing the reconstruction/utility/privacy trade-off.

\subsection{Re-identification risk}
\label{sec:defs}
Defining re-identification risk directly in terms of the URU would
measure a mechanism's privacy using the SAME component its encoder was
adversarially trained against. This paper instead defines all three
quantities below in terms of a \emph{fresh, external} attacker, trained
strictly after the encoder is frozen (Section~\ref{sec:threat-model-naive}),
for the reason Section~\ref{sec:residual-risk} demonstrates directly:
the in-model URU's own confidence stays near chance throughout
training, regardless of entropy decile, because the encoder's
adversarial training suppresses it by construction -- so a definition
of risk anchored to the URU cannot distinguish a mechanism that is
actually private from one whose own adversary is simply weak.

\textbf{Definition 1 (Re-identification risk).} Let $A$ denote the
external attacker (Section~\ref{sec:threat-model-naive}), trained
after $\mathrm{Enc}_L$ is frozen, not the URU. The re-identification
risk $R_{\mathrm{orig}}(u)$ for user $u$ is the probability that $A$
correctly identifies $u$ from their original trajectory $\tau_u$:
$R_{\mathrm{orig}}(u) = P(A(\tau_u) = u)$. The post-mechanism risk
$R_{\mathrm{anon}}(u)$ is the same attacker's accuracy on the
mechanism's released representation: $R_{\mathrm{anon}}(u) =
P(A(\mathrm{Enc}_L(\tau_u)) = u)$.

\textbf{Definition 2 (Privacy gain).} This paper reports two
decile-level statistics (Section~\ref{sec:decile-reduction-metrics}):
absolute reduction, $\overline{R_{\mathrm{orig}}}_d -
\overline{R_{\mathrm{anon}}}_d$, and relative reduction, $1 -
\overline{R_{\mathrm{anon}}}_d \,/\, \overline{R_{\mathrm{orig}}}_d$,
both computed from decile \emph{means} over users in decile $d$, with a
bootstrap confidence interval, rather than a per-user ratio
$P(\tau_u) = 1 - R_{\mathrm{anon}}(u) / R_{\mathrm{orig}}(u)$ averaged
within the decile: the per-user ratio is unstable whenever
$R_{\mathrm{orig}}(u) \approx 0$ for an individual user, since a single
such user can dominate a decile's average and flip its sign without
reflecting any real change in post-mechanism risk
(Section~\ref{sec:decile-reduction-metrics} documents a concrete case).
The per-user ratio is retained only as an appendix statistic, labeled
unstable.

\textbf{Definition 3 (Residual-risk disparity).} In place of a
percentage of users correctly re-identified within the top-$k$
predictions, this paper reports the ratio of post-mechanism risk
between the lowest- and highest-entropy deciles,
$\overline{R_{\mathrm{anon}}}_0 \,/\, \overline{R_{\mathrm{anon}}}_9$,
with a bootstrap confidence interval obtained by independently
resampling each decile's users. This statistic is reported because it
answers the paper's central question directly -- not whether the
mechanism protects users \emph{on average}, but how much worse off the
least-protected decile is than the best-protected one, after the
mechanism has been applied.

\subsection{Why these definitions}
\label{sec:methods-changed}
Three measurement choices in Definitions~1--3 are load-bearing for the
results in this paper, and are worth stating explicitly:
\begin{enumerate}
\item \textbf{A fresh external attacker, not the mechanism's own
  co-trained unit.} Every risk number in this paper comes from an
  attacker trained after the mechanism is frozen (Definition~1), never
  the co-trained URU. Section~\ref{sec:residual-risk} shows these two
  give substantially different answers to ``how private is this
  mechanism,'' because the URU's own confidence stays near chance by
  construction, independent of entropy decile, while the external
  attacker's accuracy is both higher and clearly decile-dependent.
\item \textbf{Decile-level reduction metrics with confidence intervals,
  not a single threshold crossing.} Reporting fairness only as the
  fraction of similar-user pairs whose privacy-gain disparity falls
  below a fixed threshold $\epsilon$ states whether a disparity crosses
  one, somewhat arbitrary line, not how large it is. This paper reports
  decile-level absolute and relative risk reduction and a
  residual-risk-disparity statistic (Section~\ref{sec:decile-reduction-metrics}),
  each with a bootstrap confidence interval, which states both.
\item \textbf{Multiple seeds, both attacker-side and encoder-side.}
  Every attacker-side number is averaged over five random seeds;
  results that compare across mechanism configurations additionally use
  $\geq 3$ independent encoder-training seeds, with both the mean and
  the seed-to-seed standard deviation reported (Section~\ref{sec:seeding}).
  Section~\ref{sec:phase4} shows a dominance claim that holds at 3
  encoder seeds can fail to hold at 10, because encoder-initialization
  variance alone can be large enough to manufacture an apparent
  advantage.
\end{enumerate}
